\section{SA-DMoN}
\label{sec_res_sadmon}

This section reports the results for the ten SA-DMoN configurations
introduced in Section~\ref{sec_meth_sadmon}.
The experiments begin with vanilla DMoN
and then evaluate five v1 configurations based on the skeleton-aware
spatial null model.
These vary the treatment of the spatial bandwidth $\sigma$
and include a diagnostic configuration in which the spectral term
is removed entirely.
Four v2 configurations then test whether the result changes
when the adjacency used by the modularity objective is decoupled
from the geometric message-passing graph,
when the type-exclusivity term is replaced by the entropy-based
formulation,
or when both changes are applied together.

Across these changes, the result remains largely unchanged.
All ten configurations converge to synthetic-validation PGA values
in a narrow range of approximately $0.77$--$0.79$,
below the no-depth GAT baseline of $0.8783$.
Neither changing the bandwidth treatment nor modifying the adjacency
and type-loss formulations produces a configuration that breaks out
of this range.
The consistency of the result across the two SA-DMoN generations
suggests that the limitation is not associated with a single choice
of $\sigma$ or one particular implementation of the auxiliary losses.
Instead, the experiments provide evidence that the modularity objective
itself is poorly aligned with the spatial keypoint graph used
for pose grouping.
On this basis, modularity-driven grouping is not pursued further,
and the investigation moves toward modifying the embedding network
directly.

\subsection{Setup}
\label{sec_res_sadmon_setup}

All ten SA-DMoN configurations use the standard GAT backbone of
Section~\ref{sec_meth_baseline}
and otherwise retain the optimiser and training settings of
Section~\ref{sec_meth_baseline_training}.
The experimental variables are confined to the DMoN head and,
for the v2 variants, the adjacency matrix supplied to the spectral objective.
This keeps the underlying embedding architecture fixed
while testing whether changes to the modularity formulation improve
the resulting grouping.

The v1 configurations are trained for 150 epochs.
This longer schedule was retained for variants whose head losses
were still changing at epoch 100,
so that an apparent plateau was not caused simply by terminating
training too early.
The v2 configurations are trained for 100 epochs,
after preliminary runs showed that their head losses had already
converged within this interval.
The difference in training length therefore reflects the observed
convergence behaviour of the two groups
rather than a change in the optimisation objective.

The spectral-loss coefficient $\lambda_{\text{spectral}}$ varies
between configurations and is reported with the corresponding results
below.
The orthogonality and collapse-prevention terms are both assigned
a fixed weight of $0.1$,
following the DMoN formulation of Tsitsulin et al.~\cite{tsitsulin2023dmon}.
The supervised cross-entropy term retains a weight of $1.0$ throughout.
Consequently, the ablation isolates the treatment of the modularity
signal
while keeping the remaining loss weights and GAT training configuration
fixed.

\subsection{Headline numbers}
\label{sec_res_sadmon_numbers}

\begin{table}[H]
\centering
\tablesize
\begin{tabular}{lp{0.42\linewidth}ccl}
\toprule
ID & Configuration & Synth val PGA & COCO val PGA & W\&B run \\
\midrule
& Standard GAT baseline (depth off)
    & 0.8783 & 0.8841 & \texttt{jkd1ln4a} \\
\midrule
Exp 3-bis & Vanilla DMoN, $\lambda_{\text{spectral}} = 0.1$
    & 0.7879 & 0.8657 & \texttt{124ykt2c} \\
Exp 4 & SA-DMoN v1, learnable $\sigma$, unconstrained
    & 0.7872 & 0.8658 & \texttt{69yfqshy} \\
Exp 5 & SA-DMoN v1, learnable $\sigma$, clamped $[0.05, 0.5]$
    & 0.7918 & 0.8649 & \texttt{rvd7qaow} \\
Exp 6 & SA-DMoN v1, clamped $\sigma$, $\lambda_{\text{spectral}} = 1.0$
    & 0.7892 & 0.8628 & \texttt{2ok6u2k3} \\
Exp 7 & SA-DMoN v1, $\sigma$ = per-graph median pairwise distance
    & 0.7769 & 0.8642 & \texttt{usjwc380} \\
Exp 8 & SA-DMoN v1, spectral term removed
    & 0.7889 & 0.8684 & \texttt{2ucbx3e0} \\
Exp v2a & SA-DMoN v2, feature adjacency only,
          $\lambda_{\text{spectral}} = 1.0$
    & 0.7849 & 0.8622 & \texttt{x6v3wfnv} \\
Exp v2b & SA-DMoN v2, feature adjacency only,
          $\lambda_{\text{spectral}} = 0.1$
    & 0.7925 & 0.8677 & \texttt{qn7rv98j} \\
Exp v2c & SA-DMoN v2, entropy type loss only
    & 0.7934 & 0.8660 & \texttt{zz2sj7an} \\
Exp v2d & SA-DMoN v2, feature adjacency + entropy type loss
    & 0.7910 & 0.8652 & \texttt{02jqs33r} \\
\bottomrule
\end{tabular}
\caption{Results for the ten DMoN and SA-DMoN configurations compared
with the no-depth GAT baseline. COCO val PGA is $k$-means on ground-truth
keypoints. Synthetic-validation PGA remains within
$[0.7769, 0.7934]$ across all ten configurations, a total range of
$0.0165$, and every configuration remains below the baseline value of
$0.8783$. The same pattern is observed on COCO: applying $k$-means to
the jointly trained embeddings produces PGA values between $0.8622$
and $0.8684$, compared with $0.8841$ for the baseline. The DMoN heads'
own COCO read-outs are lower still, ranging from $0.8107$ to $0.8213$.
Across the tested configurations, neither the skeleton-aware null model,
the treatment of $\sigma$, the spectral-loss weight, feature-based
adjacency, nor the entropy type loss produces an improvement over the
standard GAT baseline.}
\label{tab:sadmon_headline}
\end{table}

\subsection{Interpretation}
\label{sec_res_sadmon_discussion}

The ten configurations in Table~\ref{tab:sadmon_headline} are most
informative when considered as a sequence of attempts to change the
same underlying objective. The dominant result is not the ranking of
individual configurations, but the narrow range in which all ten
remain. Three observations follow.

\textbf{1. The plateau is the main result:}
Every DMoN and SA-DMoN configuration reaches a synthetic-validation
PGA between $0.7769$ and $0.7934$, giving a total spread of only
$0.0165$. This remains true despite substantial changes to the
modularity formulation. The spatial bandwidth $\sigma$ is made
learnable, constrained, or scene-dependent; the spectral contribution
is re-weighted or removed; the adjacency supplied to the modularity
objective is changed from geometric to feature-based; and the
type-related loss is reformulated. None of these changes moves the
result outside a band of approximately two PGA points.

More importantly, the entire band lies below the no-depth GAT baseline
of $0.8783$. The best SA-DMoN result, $0.7934$, remains $0.0849$ below
the baseline. The same pattern appears after transfer to COCO, where
$k$-means on the jointly trained embeddings remains within
$[0.8622,0.8684]$, compared with $0.8841$ for the standard GAT.
The consistency of the result across the ten configurations suggests
that the limitation is associated with the modularity-based training
signal rather than with a single poorly chosen hyperparameter.

\textbf{2. The bandwidth experiments expose a weakness in the
skeleton-aware null model:}
Experiments~4--7 vary the spatial bandwidth $\sigma$ through four
different treatments: unconstrained learning, constrained learning,
constrained learning with increased spectral weight, and a fixed
scene-adaptive value based on the median pairwise distance. None
produces a useful improvement.

The unconstrained configuration provides the clearest diagnostic.
During training, the learned bandwidth grows to approximately $341$.
From Equation~\eqref{eq:sadmon_spatial}, increasing $\sigma$ causes
the spatial kernel to approach
\begin{equation}
    P_{ij} \rightarrow 1,
\end{equation}
so the spatial component of the proposed null model progressively
loses its dependence on keypoint distance. In this limit,
Equation~\eqref{eq:sadmon_null} is governed primarily by the
joint-type affinity $T_{ab}$ rather than by the intended combination
of type and spatial information. The optimiser therefore finds a
regime in which the spatial part of the skeleton-aware null contributes
very little discrimination.

Constraining $\sigma$ to $[0.05,0.5]$ prevents this numerical growth,
but the learned value moves to the upper boundary of the allowed
interval. This behaviour is consistent with the objective favouring a
broader spatial kernel than the imposed range permits. The resulting
PGA of $0.7918$ is nevertheless only $0.0046$ above the unconstrained
configuration and remains well below the GAT baseline. Increasing
$\lambda_{\text{spectral}}$ from $0.1$ to $1.0$ reduces PGA from
$0.7918$ to $0.7892$, while replacing the learned bandwidth with the
per-scene median pairwise distance reduces it further to $0.7769$.
Within the bandwidth treatments tested here, there is therefore no
evidence of an operating point at which the skeleton-aware null
produces a competitive grouping representation.

\textbf{3. The v2 changes do not break the plateau:}
The v2 experiments test whether the v1 result is caused by two more
specific design choices. The first replaces the geometric adjacency
inside the spectral objective with an adjacency recomputed from the
current embedding, reducing the direct similarity between the
spatially constructed adjacency and spatial null. The second replaces
the original type-related term with the smoother entropy formulation
described in Section~\ref{sec_meth_sadmon}.

Feature adjacency with
$\lambda_{\text{spectral}}=1.0$ reaches $0.7849$ (v2a). Reducing the
spectral weight to $0.1$ raises this to $0.7925$ (v2b), a marginal
$0.0007$ above the best v1 configuration and far too small to alter the
overall conclusion. The entropy type loss alone gives the highest
synthetic PGA among the ten experiments at $0.7934$ (v2c), while
combining feature adjacency and entropy produces $0.7910$ (v2d).
Their COCO results remain similarly close, between $0.8622$ and
$0.8677$ for the four v2 configurations.

The v2 experiments therefore do not support either proposed change as
a solution to the SA-DMoN plateau. Decoupling the spectral adjacency
from the geometric graph and reformulating the type loss both alter the
objective, but neither produces a meaningful improvement over v1 or
approaches the standard GAT baseline. Taken together with the bandwidth
experiments, the results support the interpretation that Newman-style
modularity is poorly aligned with the spatial-proximity graph used in
this pose-grouping formulation. The graph connects joints because they
are spatially close, whereas the desired communities are defined by
person identity; proximity and community membership are therefore not
the same structural signal.

\textbf{Methodological pivot:}
The learned-head experiments of
Section~\ref{sec_res_learned_heads} and the SA-DMoN experiments above
approach grouping from substantially different directions, yet neither
produces an improvement over the standard GAT followed by $k$-means.
This does not rule out learned grouping heads in general, but it
provides little evidence that further complexity at the read-out stage
is the most productive direction under the present formulation.

The investigation therefore shifts to the representation itself.
PG-GAT, evaluated in Section~\ref{sec_res_pggat_ablations}, retains
$k$-means as the grouping read-out and instead modifies the GAT
attention mechanism to expose pose-specific structural information
directly during embedding. This change tests the hypothesis suggested
by the preceding experiments: that improving the representation
available to the grouping stage is more effective than replacing the
grouping stage itself.